\chapter{Chapter06 — Small Bias Generator}
%%% generated by the blueprint pipeline from FABL.Chapter06.Constructions.SmallBiasGenerator
\section{Overview}

Declarations of the Lean module \texttt{FABL.Chapter06.Constructions.SmallBiasGenerator}.

\section{Declarations}

\begin{definition}[Uniform pushforward density]
\lean{FABL.ProbabilityDensity.uniformPushforward}
\statementsource{Analysis-of-Boolean-Functions-by-Ryan-ODonnell}{
  pdf-ffbe9c7b3306-p152-b008,
  pdf-ffbe9c7b3306-p153-b002
}
\label{FABL.ProbabilityDensity.uniformPushforward}
\leanok
\uses{FABL.F₂Cube, FABL.ProbabilityDensity}
Let $\Omega$ be a finite nonempty set and let $g$ map $\Omega$ to the binary cube of
dimension $n$. The uniform pushforward density assigns to each cube point $x$ the value
$2^n\,|g^{-1}(x)|/|\Omega|$, where $g^{-1}(x)=\{\omega\in\Omega:g(\omega)=x\}$.
\end{definition}

\begin{theorem}[Expectation under a uniform pushforward]
\lean{FABL.ProbabilityDensity.expectation_uniformPushforward}
\label{FABL.ProbabilityDensity.expectation_uniformPushforward}
\leanok
\uses{FABL.F₂Cube, FABL.ProbabilityDensity.expectation, FABL.ProbabilityDensity.uniformPushforward}
\difficulty{5}
Let $n$ be a nonnegative integer, let $\Omega$ be a finite nonempty set, let
$g:\Omega\to\mathbb F_2^n$, and let $f:\mathbb F_2^n\to\mathbb R$. The expectation of
$f$ under the probability density obtained by pushing forward the uniform distribution
on $\Omega$ through $g$ equals the uniform expectation of $f(g(\omega))$ over
$\omega\in\Omega$.
\end{theorem}

\begin{definition}[Polynomial encoding of a binary vector]
\lean{FABL.smallBiasPolynomial}
\statementsource{Analysis-of-Boolean-Functions-by-Ryan-ODonnell}{pdf-ffbe9c7b3306-p153-b004}
\label{FABL.smallBiasPolynomial}
\leanok
\uses{FABL.BinaryExtensionField, FABL.F₂Cube, FABL.𝔽₂}
For a binary vector $\gamma=(\gamma_0,\ldots,\gamma_{n-1})$ and a binary extension field
of degree $\ell$, define the polynomial $\sum_{i=0}^{n-1}\gamma_i X^{i+1}$ over that
field, viewing each coefficient through the canonical embedding of $\mathbb{F}_2$ into
the extension field.
\end{definition}

\begin{theorem}[Evaluation of the small-bias polynomial]
\lean{FABL.smallBiasPolynomial_eval}
\label{FABL.smallBiasPolynomial_eval}
\leanok
\uses{FABL.BinaryExtensionField, FABL.F₂Cube, FABL.smallBiasPolynomial}
\difficulty{1}
Let $n$ and $[?]$ be nonnegative integers, let $[?]$ be a vector in the
$n$-dimensional binary cube, and let $r$ belong to the binary extension field of degree
$[?]$. Evaluating the polynomial whose coefficient of $X^{i+1}$ is the image of the
$i$th coordinate of $[?]$ in that field gives $\sum_{i=0}^{n-1}\gamma_i r^{i+1}$.
\end{theorem}

\begin{theorem}[Small bias polynomial ne zero]
\lean{FABL.smallBiasPolynomial_ne_zero}
\statementsource{Analysis-of-Boolean-Functions-by-Ryan-ODonnell}{pdf-ffbe9c7b3306-p153-b004}
\label{FABL.smallBiasPolynomial_ne_zero}
\leanok
\uses{FABL.BinaryExtensionField, FABL.F₂Cube, FABL.smallBiasPolynomial, FABL.𝔽₂}
\difficulty{4}
% dataset-pipeline: stub -- prose pending build_informal_statements.py
Statement pending regeneration.
\end{theorem}

\begin{theorem}[Degree bound for a binary polynomial]
\lean{FABL.smallBiasPolynomial_natDegree_le}
\proofsource{Analysis-of-Boolean-Functions-by-Ryan-ODonnell}{pdf-ffbe9c7b3306-p153-b004}
\label{FABL.smallBiasPolynomial_natDegree_le}
\leanok
\uses{FABL.F₂Cube, FABL.smallBiasPolynomial}
\difficulty{3}
For every natural number $n$ and every vector $\gamma\in(\mathbb{F}_2)^n$, form the
polynomial over the binary extension field of degree $\ell$ whose coefficient of
$X^{i+1}$ is the image of the $i$th coordinate of $\gamma$. Its natural degree is at
most $n$.
\end{theorem}

\begin{definition}[Small-bias generator]
\lean{FABL.smallBiasGenerator}
\statementsource{Vad12}{pdf-bf2ce5422088-p308-b003}
\statementsource{Analysis-of-Boolean-Functions-by-Ryan-ODonnell}{pdf-ffbe9c7b3306-p153-b003}
\label{FABL.smallBiasGenerator}
\leanok
\uses{FABL.BinaryExtensionField, FABL.F₂Cube, FABL.binaryExtensionEncode, FABL.f₂DotProduct}
Let $n$ be a natural number, let $\ell$ be a nonzero natural number, and let $r,s$ lie
in the binary extension field of degree $\ell$. The generator produces the length-$n$
binary vector whose coordinate at position $i$ is the binary dot product of the
field-coordinate encoding of $r^{i+1}$ with the encoding of $s$.
\end{definition}

\begin{theorem}[Dot product identity for the small-bias generator]
\lean{FABL.f₂DotProduct_smallBiasGenerator}
\leanok
\uses{FABL.BinaryExtensionField, FABL.F₂Cube, FABL.binaryExtensionEncode, FABL.f₂DotProduct, FABL.smallBiasGenerator, FABL.smallBiasPolynomial, FABL.smallBiasPolynomial_eval}
\difficulty{5}
Let $n$ be a natural number, let $\ell$ be a nonzero natural number, let
$\gamma\in\mathbb F_2^n$, and let $r,s$ belong to the degree-$\ell$ extension field of
$\mathbb F_2$. The dot product of $\gamma$ with the small-bias generator evaluated at
$r,s$ equals the dot product of the encoding of the polynomial with coefficients
$\gamma$ evaluated at $r$ and the encoding of $s$.
\end{theorem}

\begin{theorem}[Walsh character of the small-bias generator]
\lean{FABL.vectorWalshCharacter_smallBiasGenerator}
\label{FABL.vectorWalshCharacter_smallBiasGenerator}
\leanok
\uses{FABL.BinaryExtensionField, FABL.F₂Cube, FABL.binaryExtensionEncode, FABL.f₂DotProduct_smallBiasGenerator, FABL.smallBiasGenerator, FABL.smallBiasPolynomial, FABL.vectorWalshCharacter, FABL.vectorWalshCharacter_apply}
\difficulty{1}
Let $n$ be a natural number, let $\ell$ be a nonzero natural number, and let
$\gamma\colon \{0,1\}^n\to\mathbb F_2$, where $\mathbb F_2$ is the field with two
elements. For $r,s$ in the degree-$\ell$ extension field of $\mathbb F_2$, let
$P_\gamma(X)=\sum_{i=0}^{n-1}\gamma_iX^{i+1}$, with coefficients viewed in that
extension field, and encode its elements as binary vectors using a fixed $\mathbb
F_2$-linear isomorphism with $\mathbb F_2^\ell$. The Walsh character indexed by
$\gamma$, evaluated at the vector whose $i$th coordinate is the binary dot product of
the encodings of $r^{i+1}$ and $s$, equals the Walsh character indexed by the encoding
of $P_\gamma(r)$, evaluated at the encoding of $s$.
\end{theorem}

\begin{theorem}[Fixed-seed character expectation]
\lean{FABL.expect_smallBiasGenerator_fixed_first}
\proofsource{Analysis-of-Boolean-Functions-by-Ryan-ODonnell}{pdf-ffbe9c7b3306-p153-b004}
\label{FABL.expect_smallBiasGenerator_fixed_first}
\leanok
\uses{FABL.BinaryExtensionField, FABL.F₂Cube, FABL.binaryExtensionEncode, FABL.expect_vectorWalshCharacter, FABL.smallBiasGenerator, FABL.smallBiasPolynomial, FABL.vectorWalshCharacter, FABL.vectorWalshCharacter_smallBiasGenerator}
\difficulty{4}
Let $\ell>0$, let $\mathbb{F}$ be the degree-$\ell$ extension of $\mathbb{F}_2$, let
$\gamma\in\mathbb{F}_2^n$, and fix $r\in\mathbb{F}$. For uniform $s\in\mathbb{F}$, form
the vector $x\in\mathbb{F}_2^n$ whose $i$th coordinate is the dot product over
$\mathbb{F}_2$ of the coordinate vectors of $r^{i+1}$ and $s$, using a fixed basis of
$\mathbb{F}$ over $\mathbb{F}_2$. The mean of the Walsh character $(-1)^{\sum_i\gamma_i
x_i}$ is $1$ if $\sum_i\gamma_i r^{i+1}=0$ in $\mathbb{F}$, and is $0$ otherwise; here
each $\gamma_i$ is viewed in $\mathbb{F}$ via the natural inclusion.
\end{theorem}

\begin{theorem}[Character expectation and polynomial roots]
\lean{FABL.smallBiasGenerator_characterExpectation_eq_rootProbability}
\proofsource{Analysis-of-Boolean-Functions-by-Ryan-ODonnell}{pdf-ffbe9c7b3306-p153-b004}
\label{FABL.smallBiasGenerator_characterExpectation_eq_rootProbability}
\leanok
\uses{FABL.BinaryExtensionField, FABL.F₂Cube, FABL.binaryExtensionField_natCard, FABL.expect_smallBiasGenerator_fixed_first, FABL.smallBiasGenerator, FABL.smallBiasPolynomial, FABL.smallBiasPolynomial_ne_zero, FABL.vectorWalshCharacter}
\difficulty{5}
Let $n$ be a nonnegative integer, let $[?]$ be a positive integer, and let $[?]$ be a
nonzero vector in the $n$-dimensional binary cube. For independent uniformly sampled
elements $r,s$ of the binary extension field of degree $[?]$, the expected value of the
Walsh character indexed by $[?]$ on the generator output equals the number of roots in
that field of the polynomial associated with $[?]$, divided by $2^{[?]}$.
\end{theorem}

\begin{theorem}[Nonnegative character expectation]
\lean{FABL.smallBiasGenerator_characterExpectation_nonneg}
\proofsource{Analysis-of-Boolean-Functions-by-Ryan-ODonnell}{pdf-ffbe9c7b3306-p153-b004}
\label{FABL.smallBiasGenerator_characterExpectation_nonneg}
\leanok
\uses{FABL.BinaryExtensionField, FABL.F₂Cube, FABL.smallBiasGenerator, FABL.smallBiasGenerator_characterExpectation_eq_rootProbability, FABL.vectorWalshCharacter}
\difficulty{1}
Let $n$ be a nonnegative integer, let $[?]$ be a positive integer, and let $[?]$ be a
nonzero vector in the $n$-dimensional binary cube. The expected value of its Walsh
character on the small-bias generator output, for a uniformly sampled pair from the
binary extension field of degree $[?]$, is nonnegative.
\end{theorem}

\begin{theorem}[Upper bound on character expectation]
\lean{FABL.smallBiasGenerator_characterExpectation_le}
\proofsource{Analysis-of-Boolean-Functions-by-Ryan-ODonnell}{pdf-ffbe9c7b3306-p153-b004}
\label{FABL.smallBiasGenerator_characterExpectation_le}
\leanok
\uses{FABL.BinaryExtensionField, FABL.F₂Cube, FABL.ncard_rootSet_le_natDegree, FABL.smallBiasGenerator, FABL.smallBiasGenerator_characterExpectation_eq_rootProbability, FABL.smallBiasPolynomial, FABL.smallBiasPolynomial_natDegree_le, FABL.vectorWalshCharacter}
\difficulty{3}
Let $n$ be a nonnegative integer, let $[?]$ be a positive integer, and let $[?]$ be a
nonzero vector in the $n$-dimensional binary cube. The expected value of its Walsh
character on the small-bias generator output, for a uniformly sampled pair from the
binary extension field of degree $[?]$, is at most $n/2^{[?]}$.
\end{theorem}

\begin{definition}[Small-bias generator multiset]
\lean{FABL.smallBiasGeneratorMultiset}
\statementsource{Analysis-of-Boolean-Functions-by-Ryan-ODonnell}{pdf-ffbe9c7b3306-p153-b003}
\label{FABL.smallBiasGeneratorMultiset}
\leanok
\uses{FABL.BinaryExtensionField, FABL.F₂Cube, FABL.smallBiasGenerator}
Let $n$ be a natural number and let $\ell$ be a nonzero natural number. Take every
ordered pair $(r,s)$ of elements of the binary extension field of degree $\ell$, apply
the small-bias generator to each pair, and retain the resulting outputs as a multiset,
including repetitions.
\end{definition}

\begin{theorem}[Cardinality of the generator multiset]
\lean{FABL.smallBiasGeneratorMultiset_card}
\proofsource{Analysis-of-Boolean-Functions-by-Ryan-ODonnell}{pdf-ffbe9c7b3306-p153-b003}
\label{FABL.smallBiasGeneratorMultiset_card}
\leanok
\uses{FABL.binaryExtensionField_natCard, FABL.smallBiasGeneratorMultiset}
\difficulty{1}
Let $n$ be a nonnegative integer and let $[?]$ be a positive integer. The multiset
obtained by applying the small-bias generator to every ordered pair of elements of the
binary extension field of degree $[?]$ has cardinality $(2^{[?]})^2$.
\end{theorem}

\begin{definition}[Small bias generator density]
\lean{FABL.smallBiasGeneratorDensity}
\statementsource{Analysis-of-Boolean-Functions-by-Ryan-ODonnell}{
  pdf-ffbe9c7b3306-p152-b008,
  pdf-ffbe9c7b3306-p153-b002
}
\label{FABL.smallBiasGeneratorDensity}
\leanok
\uses{FABL.BinaryExtensionField, FABL.ProbabilityDensity, FABL.ProbabilityDensity.uniformPushforward, FABL.smallBiasGenerator}
For a natural number $n$ and a nonzero natural number $\ell$, define a probability
density on the binary cube of dimension $n$ by taking the uniform pushforward under the
small bias generator of a uniformly chosen pair of elements of the binary extension
field of degree $\ell$ over $\mathbb{F}_2$.
\end{definition}

\begin{theorem}[Fourier coefficient of a small bias generator]
\lean{FABL.vectorFourierCoeff_smallBiasGeneratorDensity}
\proofsource{Analysis-of-Boolean-Functions-by-Ryan-ODonnell}{pdf-ffbe9c7b3306-p153-b004}
\label{FABL.vectorFourierCoeff_smallBiasGeneratorDensity}
\leanok
\uses{FABL.BinaryExtensionField, FABL.F₂Cube, FABL.ProbabilityDensity.expectation, FABL.ProbabilityDensity.expectation_uniformPushforward, FABL.smallBiasGeneratorDensity, FABL.smallBiasGenerator_characterExpectation_eq_rootProbability, FABL.smallBiasPolynomial, FABL.vectorFourierCoeff, FABL.vectorFourierCoeff_eq_expect, FABL.vectorWalshCharacter}
\difficulty{3}
Let $n$ and $\ell$ be natural numbers with $\ell\ne 0$, and let
$\gamma\in\mathbb{F}_2^n$ be nonzero. The Fourier coefficient at $\gamma$ of the density
obtained by uniformly choosing two elements of the degree-$\ell$ extension field of
$\mathbb{F}_2$ and applying the small bias generator equals the number of roots in that
field of the polynomial $\sum_{i=0}^{n-1}\gamma_i X^{i+1}$, divided by $2^\ell$.
\end{theorem}

\begin{theorem}[Bias of the small-bias generator density]
\lean{FABL.smallBiasGeneratorDensity_isBiased}
\proofsource{Analysis-of-Boolean-Functions-by-Ryan-ODonnell}{
  pdf-ffbe9c7b3306-p153-b002,
  pdf-ffbe9c7b3306-p153-b004
}
\label{FABL.smallBiasGeneratorDensity_isBiased}
\leanok
\uses{FABL.ProbabilityDensity.IsBiased, FABL.smallBiasGeneratorDensity, FABL.smallBiasGenerator_characterExpectation_eq_rootProbability, FABL.smallBiasGenerator_characterExpectation_le, FABL.vectorFourierCoeff_smallBiasGeneratorDensity}
\difficulty{3}
Let $n$ and $$ be natural numbers with $\ne 0$, and let $$ be a real number
satisfying $n/2^{}\le $. The density generated by the small-bias generator over the
binary extension field of degree $$ is $$-biased: every nonzero frequency has Fourier
coefficient of absolute value at most $$.
\end{theorem}

\begin{theorem}[Small-bias generator bounds]
\lean{FABL.smallBiasGenerator_core}
\proofsource{Analysis-of-Boolean-Functions-by-Ryan-ODonnell}{
  pdf-ffbe9c7b3306-p153-b003,
  pdf-ffbe9c7b3306-p153-b004
}
\label{FABL.smallBiasGenerator_core}
\leanok
\uses{FABL.ProbabilityDensity.IsBiased, FABL.smallBiasGeneratorDensity, FABL.smallBiasGeneratorDensity_isBiased, FABL.smallBiasGeneratorMultiset, FABL.smallBiasGeneratorMultiset_card}
\difficulty{4}
Let $n,\ell$ be positive integers and let $\varepsilon>0$ satisfy $n\leq \varepsilon
2^\ell$ and $2^\ell\leq 4n/\varepsilon$. Identify the field $\mathbb{F}_{2^\ell}$ with
$\mathbb{F}_2^\ell$ by a basis, and for $r,s\in\mathbb{F}_{2^\ell}$ form the vector in
$\mathbb{F}_2^n$ whose $i$th coordinate is the dot product over $\mathbb{F}_2$ of the
coordinate vectors of $r^{i+1}$ and $s$, for $0\leq i<n$. The probability distribution
obtained by choosing $(r,s)$ uniformly has density relative to the uniform distribution
on $\mathbb{F}_2^n$ whose every nonzero Walsh coefficient has absolute value at most
$\varepsilon$; moreover, the multiset of generated vectors, counted with multiplicity,
has size at most $16(n/\varepsilon)^2$.
\end{theorem}
